\scanpage{22}
\deleted{\unsure{Lemme 3}} Alors on a $\ x_N=x_{N+1}$

\uline{Dém} Posons
\[
\left\{\begin{aligned}
&\deleted{\unsure{y_0=x_{N+1}-x_N}}\\
&y_\varepsilon=x_{N+p^\varepsilon}-x_{N\deleted{\unsure{+1}}}\qquad 0\leq\varepsilon\leq e.
\end{aligned}\right.
\]
d'où
\[
\begin{aligned}
&y_\varepsilon-y_0=x_{N+p^\varepsilon}-x_{N+1}\\
&\deleted{\unsure{y_\varepsilon=x_{N+p^\varepsilon}-x_{N+1}}}
\end{aligned}
\]
L'hypothèse s'écrit donc, en distinguant les cas $i=N,i=N+1$
\[
\begin{aligned}
&\deleted{\unsure{\binom{N+1}{N}y_0=0}}\hspace{5em}\deleted{\unsure{(i=N,j=N+1)}}\\
&(\deleted{\unsure{\bar S}}S)=S(N,e,(y_\varepsilon)_{0\leq\varepsilon\leq e})\left\{\begin{aligned}
&c_{N,p^\varepsilon}y_\varepsilon=0\qquad 0\leq\varepsilon\leq e\\
&c_{N+1,p^\varepsilon-1}(y_\varepsilon-y_0)=0\qquad 0\leq\varepsilon\leq e\quad\text{(le cas $\varepsilon=0$ étant trivial)}
\end{aligned}\right.
\end{aligned}
\]
La conclusion s'écrit $y_0=0$. En fait, nous allons déduire cette conclusion du système plus faible \inserted{(équivalent si $E$ est \deleted{un} annulé par $p$, i.e. est un $\mathbb{F}_p$-vectoriel)}
\[
(\bar S^{\unsure{\prime|{}}})=\bar S(N,e,(y_\varepsilon)_{0\leq\varepsilon\leq e}\left\{\begin{aligned}
&y_\varepsilon=0\quad\text{si }c_{N,p^\varepsilon}\not\equiv0\ (p)\qquad 0\leq\varepsilon\leq e\\
&y_\varepsilon-y_0=0\quad\text{si }c_{N+1,p^\varepsilon-1}\not\equiv0\ (p)\qquad 1\leq\varepsilon\leq e.
\end{aligned}\right.
\]

Pour prouver que ce système implique $y_0=0$, on récurre sur $N$. Comme la première équation pour $\varepsilon=0$ donne
\[
y_0=0\quad\text{si}\quad N+1\not\equiv0\ (p)
\]
la conclusion est triviale pour $N=0$. \deleted{\unsure{Alors on}} On suppose donc $N>0$, \inserted{et prouvé pour les $N'<N$.} La conclusion est triviale si $N+1\not\equiv0\ (p)$. Donc OPS
\[
N+1\equiv0\ (p)\qquad\text{soit}\qquad N+1=(N'+1)p
\]
où $N'=\frac{N+1}{p}-1<N$ (car $\frac{N+1}{p}<N+1$). Alors le cor. précédent montre que le syst. $\bar S=\bar S(N,e,(y_\varepsilon)_{0\leq\varepsilon\leq e})$ \inserted{(compte tenu de $N+1\equiv0\ (p)$ qui nous permet d'oublier la première équation pour $\varepsilon=0$)} équivaut au système (où $e'=e-1$, $y'_\varepsilon=y_{\varepsilon+1}$ pour $0\leq\varepsilon\leq e'=e-1$
\[
\deleted{\unsure{\bar S'=\bar S'(N',e',(y'_\varepsilon)_{1\leq\varepsilon}}}\quad
\left\{\begin{aligned}
&y'_\varepsilon=0\quad\text{si }c_{N',p^\varepsilon}\not\equiv0\ (p)\qquad 0\leq\varepsilon\unsure{=|\leq}e'\\
&y'_\varepsilon-y_0=0\quad\text{si }c_{N'+1,p^\varepsilon-1}\not\equiv0\ (p)\qquad 0\leq\varepsilon\leq e'
\end{aligned}\right.
\]
Pour $\varepsilon=0$ cette dernière équation donne \deleted{\unsure{$x$}}
\[
y'_0=y_0
\]
donc on peut se borner à $1\leq\varepsilon\leq e'$ (et se rappeler $y'_0=y_0$).
% archive sheet number: 21

